\documentclass[11pt,reqno]{amsart}
\usepackage[T1]{fontenc}
\usepackage{lmodern}
\usepackage{amsmath,amssymb,mathtools}
\usepackage{microtype}
\usepackage{needspace}
\usepackage[margin=1.05in]{geometry}
\usepackage[hidelinks,unicode]{hyperref}
\urlstyle{same}
\hypersetup{pdftitle={Maximum-modulus points on concentric circles},
 pdfauthor={Yunpeng Li},
 pdfsubject={A locally finite exceptional-set bound for entire functions},
 pdfkeywords={entire function; maximum modulus; analytic curve; multiplicative correspondence}}

\newtheorem{theorem}{Theorem}[section]
\newtheorem{lemma}[theorem]{Lemma}
\newtheorem{proposition}[theorem]{Proposition}
\newtheorem{corollary}[theorem]{Corollary}
\theoremstyle{remark}
\newtheorem{remark}[theorem]{Remark}
\newtheorem{example}[theorem]{Example}
\numberwithin{equation}{section}
\newcommand{\C}{\mathbb C}
\newcommand{\R}{\mathbb R}
\newcommand{\N}{\mathbb N}
\newcommand{\Cs}{\C^*}
\newcommand{\Ps}{\widehat{\C}}
\newcommand{\ord}{\operatorname{ord}}
\newcommand{\reg}{\mathrm{reg}}
\newcommand{\sing}{\mathrm{sing}}
\newcommand{\red}{\mathrm{red}}
\DeclareMathOperator{\Rea}{Re}
\DeclareMathOperator{\Ima}{Im}
\DeclareMathOperator{\Res}{Res}
\allowdisplaybreaks[1]

\title[Maximum-modulus points on circles]{Maximum-modulus points on concentric circles}
\author{Yunpeng Li}
\address{School of Data Science, The Chinese University of Hong Kong,
Shenzhen, China}
\email{liyunpeng@cuhk.edu.cn}
\date{October 11, 2026}
\subjclass[2020]{Primary 30D15; Secondary 30D20, 32B10}
\keywords{Entire function, maximum modulus, analytic curve, finite holomorphic map, multiplicative correspondence}

\begin{document}
\begin{abstract}
Let $f$ be a nonzero, non-monomial entire function, and let $\nu_f(r)$ count the
points at which $f$ attains its maximum modulus on $|z|=r$.
Write $f(z)=z^m(a_0+a_kz^k+O(z^{k+1}))$, where $a_0a_k\ne0$.
We prove that $\nu_f(r)\le 2k$ outside a locally finite subset of
$(0,\infty)$. The constant $2k$ is sharp. In particular,
$\liminf_{r\to\infty}\nu_f(r)<\infty$, so the number of maximum-modulus
points of a fixed non-monomial entire function cannot tend to infinity
on all sufficiently large circles. The proof associates the logarithmic
derivative $q=zf'/f$ with its reflection $q^*(w)=\overline{q(\bar w)}$
through the correspondence $q(z)=q^*(w)$ and the map
$(z,w)\mapsto(zw,q(z))$.
Local valence at the origin bounds the number of inverse pairs that can
escape. A local elimination argument then bounds the generic degree of
each nonvertical image component. The residue structure of $f'/f$
excludes vertical components.
\end{abstract}
\maketitle

\section{Introduction}
For a nonzero entire function $f$, let
\[
 M_f(r)=\max_{|z|=r}|f(z)|,
 \qquad
 \nu_f(r)=\#\{z:|z|=r,\ |f(z)|=M_f(r)\},\qquad r>0.
\]
If $f(z)=cz^m$, with $c\ne0$ and $m\ge0$, every point of each circle is
a maximum-modulus point. Throughout this paper, a non-monomial entire
function means a nonzero entire function that is not of this form.
For such a function, $\nu_f(r)$ is finite for every $r>0$, as proved
in Lemma~\ref{lem:finite-maxima}.

A question of Erd\H{o}s, recorded in Hayman's problem collection
\cite{Hayman1967}, distinguishes two possibilities for a single
non-monomial entire function: unbounded values of $\nu_f$ along a
sequence of radii, and divergence of $\nu_f(r)$ as $r\to\infty$.
The all-radii version has also been attributed to Clunie; see the
historical discussion in \cite{GlucksamPardo2024}.
Herzog and Piranian \cite{HerzogPiranian1968} constructed an entire
function with $\nu_f(n)=n$ for every positive integer $n$.
Pardo-Sim\'on and Sixsmith \cite{PardoSixsmith2026} subsequently obtained
unbounded selected-radius multiplicity for a finite-order entire
function in the Eremenko--Lyubich class $\mathcal B$.
Neither assertion controls the intervening radii.

Gl\"ucksam and Pardo-Sim\'on \cite{GlucksamPardo2024} proved an approximate
all-radii result: for every fixed $\varepsilon>0$, the number of
components of
\[
 \{z:|z|=r,\ |f(z)|>M_f(r)-\varepsilon\}
\]
tends to infinity for one entire function. The components need not all
contain exact maximum-modulus points. Marchenko
\cite[p.~215, Conjecture~1]{Marchenko2012} conjectured that
$\liminf_{r\to\infty}\nu_f(r)<\infty$ for every non-monomial entire
function of finite lower order. We prove this conclusion without a
growth assumption, together with a sharp quantitative bound.

\Needspace{9\baselineskip}
\begin{theorem}\label{thm:main}
Let $f$ be a non-monomial entire function. Write
\begin{equation}\label{eq:firstgap}
 f(z)=z^m\bigl(a_0+a_kz^k+O(z^{k+1})\bigr),
 \qquad m\ge0,\quad k\ge1,\quad a_0a_k\ne0.
\end{equation}
Then the set
\[
 E_f=\{r>0:\nu_f(r)>2k\}
\]
is locally finite in $(0,\infty)$. The constant $2k$ is sharp.
\end{theorem}

Here $m=\ord_0 f$ and $k=\ord_0(g-g(0))$, where $g=z^{-m}f$.
Thus $k$ is the first positive exponent occurring in $g$.
This definition differs from the inner degree, which is the greatest
common divisor of the positive exponents occurring in $g$; the two
integers can coincide. See \cite{EvdoridouPardoSixsmith2021}.
Local finiteness means that $E_f$ meets every compact subinterval of
$(0,\infty)$ in finitely many points.

\begin{corollary}\label{cor:negative}
For every non-monomial entire function satisfying \eqref{eq:firstgap},
\[
 \liminf_{r\to\infty}\nu_f(r)\le2k.
\]
In particular, no non-monomial entire function satisfies
$\nu_f(r)\to\infty$ as $r\to\infty$.
\end{corollary}

\begin{proof}
Each interval $(n,n+1)$ contains a radius outside $E_f$. Choosing one
such radius for every positive integer $n$ proves the assertion.
\end{proof}

Local finiteness permits an unbounded sequence of exceptional radii.
Thus Theorem~\ref{thm:main} is consistent with the selected-radius
constructions above. Section~\ref{sec:examples} gives elementary
examples showing both sharpness of the bound and the necessity of
allowing exceptional radii.

\subsection*{Outline of the proof}
Set $q=zf'/f$ and $q^*(w)=\overline{q(\bar w)}$.
At a differentiability point of $H(x)=\log M_f(e^x)$, all maxima on the
circle $|z|=e^x$ satisfy $q(z)=H'(x)\in\R$. They therefore determine
inverse pairs of
\[
 (z,w)\longmapsto (zw,q(z)),\qquad q(z)=q^*(w),
\]
over the single point $(e^{2x},H'(x))$.
This map is proper away from the value $q(0)$ in its second coordinate.
Its image is a complex-analytic curve, constructed locally from finite
products of inverse roots.

A component whose generic fiber has more than $2k$ pairs cannot lose
all its pairs through zero or infinity as the second coordinate tends
to $q(0)$. At least one bounded pair survives. This supplies a local
polynomial equation across the omitted value, and the image component
extends analytically. Its projection to the product coordinate is
proper. Unless the component is vertical, it must reach arbitrarily
small nonzero products, where local valence gives at most $2k$ pairs.
Vertical components are incompatible with the positive residues of
$f'/f$. This proves the bound outside a null set of radii. A local
parametrization of stationary arcs then makes the exceptional set
locally finite.

\section{A multiplicative correspondence}\label{sec:correspondence}
We first separate the analytic argument from the entire-function
application. Let $A$ and $B$ be nonconstant meromorphic functions on
$\C$, holomorphic at zero, such that
\begin{equation}\label{eq:generalorigin}
 A(0)=B(0)=c\in\C,
 \qquad k_A=\ord_0(A-c),\quad k_B=\ord_0(B-c).
\end{equation}
Let $L=k_A+k_B$, and write $\Cs=\C\setminus\{0\}$.
All equalities of meromorphic functions below are understood as
equalities of maps to the Riemann sphere $\Ps$.
An inverse pair over $(s,p)$ means an ordered pair $(z,w)\in(\Cs)^2$
with $zw=s$ and $A(z)=B(w)=p$.
All fiber cardinalities count distinct inverse pairs; root
multiplicities enter only in the local polynomials and mapping degrees.

Consider the reduced analytic curve
\[
 \mathcal C=\{(z,w)\in(\Cs)^2:A(z)=B(w)\},
\]
and the holomorphic map
\begin{equation}\label{eq:Phi}
 \Phi:\mathcal C\longrightarrow\Cs\times\Ps,
 \qquad \Phi(z,w)=(zw,A(z)).
\end{equation}
Put
\[
 \Omega=\Cs\times(\Ps\setminus\{c\}),
 \qquad \mathcal C^\circ=\Phi^{-1}(\Omega),
 \qquad Y=\Phi(\mathcal C^\circ).
\]
The source has pure complex dimension one. Indeed, near a pair
$(z_0,w_0)$ with finite common value $p_0$, write
\[
 A(z)-p_0=(z-z_0)^a h_A(z),\qquad
 B(w)-p_0=(w-w_0)^b h_B(w),
\]
where $a,b\ge1$ and the two holomorphic factors are nonvanishing.
Choosing local holomorphic roots of these factors gives conformal
coordinates
\[
 u=(z-z_0)h_A(z)^{1/a},\qquad
 v=(w-w_0)h_B(w)^{1/b},
\]
in which the correspondence is $u^a=v^b$. If $d=\gcd(a,b)$, its
$d$ local branches admit parametrizations
\[
 u=\zeta t^{b/d},\qquad v=t^{a/d},\qquad
 \zeta^{a/d}=\omega,\quad \omega^d=1.
\]
One choice of $\zeta$ for each $\omega$ suffices. These are
one-to-one parametrizations of the respective curve germs, because
$a/d$ and $b/d$ are coprime. At a common pole, apply the same
argument to $1/A$ and $1/B$, whose common value is zero.

\begin{lemma}[Small-coordinate fibers]\label{lem:small}
There exists $\varepsilon>0$ such that, for every
$(s,p)\in\Cs\times\Ps$, at most $L$
pairs satisfying
\begin{equation}\label{eq:pairs}
 zw=s,\qquad A(z)=B(w)=p
\end{equation}
have $\min(|z|,|w|)<\varepsilon$. In particular, if
$0<|s|<\varepsilon^2$, the whole fiber has at most $L$ distinct pairs.
We may also choose $\varepsilon$ so that neither $A$ nor $B$ takes the
value $c$ in $0<|z|<\varepsilon$.
\end{lemma}

\begin{proof}
Write $A(z)-c=z^{k_A}h_A(z)$ with $h_A(0)\ne0$, and choose a
holomorphic $k_A$th root of $h_A$ in a small disk. Then
\[
 \tau_A(z)=z h_A(z)^{1/k_A},\qquad
 \tau_A'(0)=h_A(0)^{1/k_A}\ne0,\qquad
 A(z)-c=\tau_A(z)^{k_A}.
\]
Shrink the disk until $\tau_A$ is injective and $A$ has no poles
there. For finite $p$, the equation $\tau_A(z)^{k_A}=p-c$ has at
most $k_A$ distinct solutions, even if some $k_A$th roots of $p-c$
lie outside the coordinate image. For $p=\infty$ there are no
solutions in the disk. Repeat this for $B$, and take a common disk
$|z|<\varepsilon$ on which both conclusions hold.
For fixed $s\ne0$, choosing either coordinate determines the other.
There are therefore at most $k_A$ pairs with small first coordinate and
at most $k_B$ with small second coordinate. If
$|s|<\varepsilon^2$, every pair has a small coordinate. The assertion
about the value $c$ follows from the same local coordinates.
\end{proof}

\subsection{Local product equations}
The elementary elimination fact used in both the image and extension
arguments is stated explicitly.

\begin{lemma}[Local product polynomial]\label{lem:product}
Let $A$ and $B$ be holomorphic on neighborhoods of finite unions of
closed, pairwise disjoint disks $U_i$ and $V_j$, respectively. Suppose
$A-p_0$ and $B-p_0$ have no zeros on the disk boundaries. List the roots
of $A(z)=p$ in the $U_i$ as $\alpha_1(p),\ldots,\alpha_N(p)$, and the
roots of $B(w)=p$ in the $V_j$ as $\beta_1(p),\ldots,\beta_M(p)$,
with multiplicity. For $p$ near $p_0$,
\begin{equation}\label{eq:productpolynomial}
 P(s,p)=\prod_{i=1}^N\prod_{j=1}^M
       \bigl(s-\alpha_i(p)\beta_j(p)\bigr)
\end{equation}
is a monic polynomial in $s$ with holomorphic coefficients in $p$.
Its zero set consists exactly of the products of these local inverse
roots. In particular, $P(s,p_0)$ is not identically zero as a
polynomial in $s$.
\end{lemma}

\begin{proof}
If either union of disks is empty, the assertion holds with $P=1$.
Assume that both unions are nonempty.
The minimum of $|A-p_0|$ over the finitely many boundaries
$\partial U_i$ is positive, and similarly for $B-p_0$ on the
$\partial V_j$. Choose a parameter disk smaller than both minima.
Rouch\'e's theorem makes the total root counts $N$ and $M$ constant
on that disk. If either count is zero, then again $P=1$; hence assume
$NM\ge1$. For $n\ge1$, the argument principle gives holomorphic
power sums
\[
 u_n(p)=\frac{1}{2\pi i}\sum_i
 \int_{\partial U_i}\frac{z^n A'(z)}{A(z)-p}\,dz,
 \qquad
 v_n(p)=\frac{1}{2\pi i}\sum_j
 \int_{\partial V_j}\frac{w^n B'(w)}{B(w)-p}\,dw.
\]
The denominators stay nonzero on the fixed contours, so the
integrals depend holomorphically on $p$. The residue at a root
$\alpha$ of multiplicity $r$ is $r\alpha^n$. Consequently,
\[
 u_n(p)=\sum_{i=1}^N\alpha_i(p)^n,\qquad
 v_n(p)=\sum_{j=1}^M\beta_j(p)^n,
\]
without any need to choose single-valued inverse branches. The
$n$th power sum of the $NM$ product roots is
\[
 T_n(p)=\sum_{i=1}^N\sum_{j=1}^M
          (\alpha_i(p)\beta_j(p))^n=u_n(p)v_n(p).
\]
Put $n_P=NM$ and write
\[
 P(s,p)=s^{n_P}+c_1(p)s^{n_P-1}+\cdots+c_{n_P}(p),\qquad c_0(p)=1.
\]
Newton's identities give the explicit recursion
\begin{equation}\label{eq:newtonrecursion}
 \ell c_\ell(p)+
 \sum_{n=1}^{\ell}c_{\ell-n}(p)T_n(p)=0,
 \qquad 1\le\ell\le n_P.
\end{equation}
Induction on $\ell$ proves that every $c_\ell$ is holomorphic.
For example, $c_1=-T_1$ when $n_P\ge1$, and
$c_2=(T_1^2-T_2)/2$ when $n_P\ge2$.
At each fixed $p$, the same identities recover the coefficients of
the actual finite product in \eqref{eq:productpolynomial}. Thus its
roots are exactly the products of the local inverse roots, including
at coalescence values. At $p=p_0$ the polynomial specializes to
\[
 P(s,p_0)=\prod_{\alpha,\beta}
 (s-\alpha\beta)^{\ord_\alpha(A-p_0)\,\ord_\beta(B-p_0)},
\]
where the product runs over the central roots in the chosen disks.
This polynomial is monic, so it cannot vanish identically in $s$.
\end{proof}

Individual inverse roots in Lemma~\ref{lem:product} may ramify; no
single-valued labeling is asserted. If one collection is empty,
the empty product is $P=1$. Near a common pole use the holomorphic
functions $1/A$, $1/B$ and the parameter $\lambda=1/p$. Their local
root products give $P(s,\lambda)$ by exactly the same recursion.

\begin{lemma}[Properness and the analytic image]\label{lem:properimage}
The map $\Phi:\mathcal C^\circ\to\Omega$ is proper with finite
fibers. Its image $Y$, with the reduced structure, is a closed pure
one-dimensional analytic subset of $\Omega$.
\end{lemma}

\begin{proof}
Let $K\subset\Omega$ be compact. Its $p$-projection is a compact
subset of $\Ps$ avoiding $c$. Continuity of $A$ and $B$ at zero gives
$\eta>0$ such that every inverse pair over $K$ has
$|z|,|w|\ge\eta$. If $|s|\le S$ on $K$, then
\begin{equation}\label{eq:properbounds}
 \eta\le|z|,|w|\le S/\eta.
\end{equation}
The inverse image of $K$ is closed in this compact annulus product,
and is therefore compact. For a fixed $(s,p)$, the equation $A(z)=p$
has finitely many roots in that annulus, and $w=s/z$ is determined.
This includes $p=\infty$.

Properness makes $Y$ closed. To obtain its local equations, fix
$y_0=(s_0,p_0)\in Y$ and take small disjoint disks around every first
and second coordinate occurring in its finite fiber. For finite $p_0$
the disks avoid poles; for $p_0=\infty$ use reciprocal value
coordinates. Properness implies that every inverse pair over a
sufficiently small target neighborhood lies in the products of these
disks. Otherwise a sequence of inverse pairs outside them would have
a subsequence converging to the central fiber.

Apply Lemma~\ref{lem:product} to the coordinate disks. Every nearby
point of $Y$ satisfies $P(s,p)=0$. Conversely, each nearby zero of $P$
is the product of actual inverse roots, so it belongs to $Y$. Thus
$Y=\{P=0\}_{\red}$ locally. The monic polynomial is nonzero and gives
a hypersurface in a two-dimensional coordinate neighborhood, hence a
pure analytic curve. No source component is contracted, since all
fibers are finite.
\end{proof}

\subsection{Generic degrees}
We use the standard local structure of reduced analytic curves: finite
irreducible decomposition, one-to-one disk parametrizations of
irreducible curve germs, connected punctured branches, and a locally
finite singular set. The regular locus of a globally irreducible
curve is connected. These facts, and the local finite-covering
statement, can be found in
\cite[\S6.7: Proposition~6.7.3, Theorems~6.7.4 and~6.7.6,
and Exercise~6.7.3]{Lebl2026}.
For clarity, a normalization separates the finitely many branches
at a singular curve point: each local branch is represented by its
own holomorphic disk parameter. The resulting Riemann surface is
connected for an irreducible curve.

For an irreducible component $\Gamma$ of $Y$, define $d_\Gamma$ to be
the generic number of distinct pairs in a fiber over $\Gamma$,
including all source components above it.

\begin{lemma}[Generic degree and smooth fibers]\label{lem:degree}
The integer $d_\Gamma$ is finite and well defined. If $y\in\Gamma$ is
a smooth point of the total reduced image $Y$, then
\begin{equation}\label{eq:smoothfiber}
 \#\Phi^{-1}(y)\le d_\Gamma.
\end{equation}
Generic points with exactly $d_\Gamma$ distinct inverse pairs are
dense in $\Gamma$.
\end{lemma}

\begin{proof}
Fix a smooth point $y_0\in\Gamma$ of the total image $Y$, and choose a
holomorphic coordinate $u$ on a target disk in $Y$, with $u(y_0)=0$.
The inverse fiber is finite. Parametrize every local source branch
above its points by a disk map $\gamma_\ell(t)$, with
$\gamma_\ell(0)$ in that fiber. Each parametrization is one-to-one
away from its central singular point. In a sufficiently small
neighborhood the only image branch is the target disk. Hence
\[
 h_\ell(t)=u\bigl(\Phi(\gamma_\ell(t))\bigr)
\]
is holomorphic. It is nonconstant: otherwise a whole source branch
would lie in one fiber, contrary to Lemma~\ref{lem:properimage}.
Write
\[
 h_\ell(t)=a_\ell t^{e_\ell}(1+t b_\ell(t)),
 \qquad a_\ell\ne0,\quad e_\ell\ge1.
\]
The change of coordinate
\[
 v_\ell=t\bigl[a_\ell(1+t b_\ell(t))\bigr]^{1/e_\ell}
\]
is conformal near zero and gives $h_\ell=v_\ell^{e_\ell}$.
Thus this normalized source branch contributes exactly $e_\ell$
preimages to each sufficiently small nonzero target value, and
one normalized preimage to the central value.

Choose disjoint neighborhoods of the distinct source points and
small representatives of all their branches. Properness ensures
that every inverse pair over a sufficiently small target disk is
in these neighborhoods. Indeed, omitted pairs over a sequence
tending to $y_0$ would have a convergent subsequence in the
central fiber. No inverse pairs are therefore missing from the
local count. For generic nearby values the branch preimages are
also distinct as source points: two distinct local branches
intersect only at their center. The total local degree is
\begin{equation}\label{eq:localsumdegree}
 D(y_0)=\sum_\ell e_\ell.
\end{equation}
At any point of this target disk the number of normalized
preimages, counted without multiplicity, is at most $D(y_0)$.
The number of actual source points is no larger, since different
normalized points can represent the same singular source point.
This proves the local form of \eqref{eq:smoothfiber}.

The degrees in \eqref{eq:localsumdegree} agree on overlapping
target disks: both equal the actual fiber cardinality at a generic
point of their overlap. The regular locus of $\Gamma$ is
connected. Deleting its intersections with other image components
deletes a locally finite set; a connected Riemann surface remains
connected after such a deletion. The local degrees therefore give
one finite positive integer $d_\Gamma$ on all smooth points of the
total image belonging to $\Gamma$.

The same local normal forms show that, after shrinking a smooth
target disk centered at $y_0$, every point other than its center
has exactly $D(y_0)=d_\Gamma$ distinct inverse pairs. Thus the
exceptional fibers form a locally finite set within the smooth
locus of $Y$. The singular locus of $Y$ is also locally finite,
and each branch at a singular point has a punctured smooth
representative. Generic points are therefore dense in the whole
of $\Gamma$, as asserted.
\end{proof}

The smoothness requirement in \eqref{eq:smoothfiber} concerns the
\emph{total} image $Y$, not merely one component. At a singular image
point, pairs belonging to several local image branches may occur in
the same fiber.
Also, $d_\Gamma$ is a degree of the source-to-image map. The degree
$NM$ of a local product polynomial in the variable $s$ counts
products of local inverse roots, with multiplicity. These are
different quantities, and no equality between them is used.

\subsection{Continuation at the omitted value}
The following lemma is the main continuation step.

\begin{lemma}[Continuation across the omitted value]\label{lem:closure}
If $d_\Gamma>L$, then the closure of $\Gamma$ in
$\Cs\times\Ps$ is a closed irreducible analytic curve. Moreover, if
\[
 \Sigma=\{\alpha\beta:\alpha,\beta\ne0,
                    \ A(\alpha)=B(\beta)=c\},
\]
then
\begin{equation}\label{eq:boundary}
 \overline\Gamma\cap(\Cs\times\{c\})
 \subset\Sigma\times\{c\}.
\end{equation}
The set $\Sigma$ is locally finite and contains no $s$ with
$0<|s|<\varepsilon^2$.
\end{lemma}

\begin{proof}
The roots of $A-c$ and $B-c$ in $\Cs$ avoid the disk
$|z|<\varepsilon$, by Lemma~\ref{lem:small}. A bound
$|\alpha\beta|\le S$ therefore bounds both root coordinates between
$\varepsilon$ and $S/\varepsilon$. Only finitely many roots occur in
that compact annulus. This proves the assertions about $\Sigma$.

Fix a boundary point $y_0=(s_0,c)$ of $\Gamma$, with $s_0\ne0$, and
choose a target neighborhood on which $|s|\le S$. Every generic point
of $\Gamma$ in this neighborhood has more than $L$ distinct inverse
pairs. At most $L$ have a coordinate smaller than $\varepsilon$.
Consequently, at least one pair satisfies
\begin{equation}\label{eq:survivingpair}
 \varepsilon\le |z|,|w|\le S/\varepsilon.
\end{equation}
A sequence of such pairs over generic points tending to $y_0$ has a
convergent subsequence. Its limit $(\alpha,\beta)$ satisfies
$A(\alpha)=B(\beta)=c$ and $\alpha\beta=s_0$.
This proves \eqref{eq:boundary}. The chosen pair may change from one
target point to the next.

We next obtain one analytic equation containing the whole nearby
component. There are finitely many roots of $A-c$ and $B-c$ in an
enlarged compact annulus containing \eqref{eq:survivingpair}. Choose
small coordinate disks around all of them. For $p$ sufficiently close
to $c$, every pair satisfying \eqref{eq:survivingpair} has its
coordinates in those disks. This follows from compactness and
continuity; a contrary sequence would converge to an omitted root
of $A-c$ or $B-c$.

Lemma~\ref{lem:product} now supplies a polynomial $P_0(s,p)$, monic
in $s$ and holomorphic in $p$ across $c$, from the inverse roots in
these disks. Every generic point of $\Gamma$ sufficiently near $y_0$
satisfies $P_0(s,p)=0$, because one bounded inverse pair is available.
Generic points are dense, and $P_0$ is continuous, so there is a
neighborhood $U$ of $y_0$ for which
\begin{equation}\label{eq:componentcontainment}
 \Gamma\cap U\subset
 Z:=\{(s,p)\in U:P_0(s,p)=0\}_{\red}.
\end{equation}
Only one bounded inverse pair per generic target point was needed
for this containment; the other inverse pairs may escape. The
central specialization is
\begin{equation}\label{eq:centralP}
 P_0(s,c)=\prod_{\alpha,\beta}
 (s-\alpha\beta)^{\ord_\alpha(A-c)\,\ord_\beta(B-c)},
\end{equation}
which is not identically zero.

Since $P_0(s,c)$ is a nonzero polynomial, choose a small $s$-disk
around $s_0$ in which $s_0$ is its only root and whose boundary
contains no root. After shrinking the $p$-disk around $c$, $P_0$
is still nonzero on that $s$-boundary. The local curve $Z$ thus
has finitely many irreducible branches at $y_0$, none contained
in $p=c$. For each branch choose a one-to-one normalized disk
parametrization $\gamma_j(t)$ centered at $y_0$. Its $p$-coordinate
satisfies
\[
 p(\gamma_j(t))-c=a_jt^{n_j}(1+O(t)),\qquad a_j\ne0.
\]
A holomorphic change of parameter reduces this to
$p(\gamma_j(t))-c=t^{n_j}$. In particular, the branch meets
$p=c$ only at its center, and its punctured representative is
connected. Shrink the representatives further so that they are
smooth and pairwise disjoint away from $y_0$.

On each punctured branch $Z_j^*$, the set $\Gamma\cap Z_j^*$
is closed, because $\Gamma$ is closed in $\Omega$ and
$Z_j^*\subset\Omega$. It is also open. To verify this, let
$y\in\Gamma\cap Z_j^*$. By \eqref{eq:componentcontainment},
the germ of $\Gamma$ at $y$ is a one-dimensional analytic subset
of the smooth one-dimensional germ of $Z_j^*$. In a disk
coordinate, a proper analytic subset is discrete; it cannot
have dimension one. Hence $\Gamma$ contains a neighborhood of
$y$ in $Z_j^*$. Connectedness now gives
\[
 \Gamma\cap Z_j^*=Z_j^*\quad\hbox{or}\quad
 \Gamma\cap Z_j^*=\varnothing.
\]

It follows that the closure of $\Gamma$ near $y_0$ is a union of a
subset of the finitely many full local branches of $Z$, and is
analytic. At points of $p=c$ not in the closure there is nothing to
extend. This gives a closed analytic curve globally. It remains
irreducible because $\Gamma$ is dense and irreducible. Indeed, a
decomposition of the closure into two proper closed analytic subsets
would restrict to such a decomposition of $\Gamma$: neither subset
could contain all of $\Gamma$, since it would then contain its
closure. No curve component is added inside
$p=c$, by \eqref{eq:centralP}.
\end{proof}

\begin{theorem}[Degree bound]\label{thm:correspondence}
For every irreducible component $\Gamma$ of $Y$ on which the product
coordinate $s$ is nonconstant,
\[
 d_\Gamma\le k_A+k_B.
\]
\end{theorem}

\begin{proof}
Suppose $d_\Gamma>L$. By Lemma~\ref{lem:closure}, the analytic closure
$\overline\Gamma$ lies in $\Cs\times\Ps$. Its projection
\[
 \pi:\overline\Gamma\longrightarrow\Cs,\qquad \pi(s,p)=s,
\]
is proper: for compact $K\subset\Cs$, its inverse image is a closed
subset of the compact set $K\times\Ps$. In particular, its image is
closed.

The projection is also open. Pull $s$ back to the connected
normalization of the irreducible curve $\overline\Gamma$.
It is a holomorphic function there. If $s$ were constant on one
local branch, this pullback would be constant on an open disk,
and hence constant everywhere by the identity theorem. That
would make $s$ constant on $\Gamma$, contrary to the hypothesis.
Thus $s$ is nonconstant on every parametrizing disk. The open
mapping theorem applies at its center even when its derivative
vanishes, and shows that the image of every neighborhood of
every point contains a neighborhood of that point's $s$-value.
Thus $\pi(\overline\Gamma)$ is nonempty, open, and closed in
$\Cs$, and equals $\Cs$.

In particular, $\overline\Gamma$ has points with
$0<|s|<\varepsilon^2$. Such points cannot have $p=c$, by
Lemma~\ref{lem:closure}, and therefore lie in $\Gamma$. This nonempty
open part of $\Gamma$ contains generic points with $d_\Gamma>L$
inverse pairs. Lemma~\ref{lem:small} bounds every fiber there by $L$,
a contradiction.
\end{proof}

\begin{remark}\label{rem:vertical}
The nonconstant-product hypothesis is necessary for general
meromorphic functions. For
\[
 A(z)=B(z)=\frac{z+z^3}{1+z^4},
\]
one has $A(0)=0$, $A'(0)=1$, and $A(z)=A(1/z)$.
The image component $s=1$ has four inverse pairs over a generic $p$,
although $k_A+k_B=2$. The entire-function application excludes these
vertical components by a residue argument.
Indeed, $A(z)=p$ is equivalent to
\[
 pz^4-z^3-z+p=0.
\]
For a generic nonzero $p$ this quartic has four distinct nonzero
roots, producing four distinct pairs $(z,1/z)$ over $(1,p)$.
\end{remark}

\section{Maximum-modulus points}\label{sec:application}
Let $f$ satisfy \eqref{eq:firstgap}, and set
\begin{equation}\label{eq:q}
 q(z)=\frac{zf'(z)}{f(z)},\qquad
 q^*(w)=\overline{q(\bar w)}.
\end{equation}
To compute the expansion explicitly, write $f=z^m g$. Then
\[
 q(z)=m+\frac{zg'(z)}{g(z)},\qquad
 g(z)=a_0+a_kz^k+O(z^{k+1}),\qquad
 zg'(z)=ka_kz^k+O(z^{k+1}).
\]
Since $1/g(z)=a_0^{-1}+O(z^k)$, the removable value at zero
and the first nonconstant term are
\begin{equation}\label{eq:qexpansion}
 q(0)=q^*(0)=m,\qquad
 q(z)=m+\frac{k a_k}{a_0}z^k+O(z^{k+1}).
\end{equation}
Apply Section~\ref{sec:correspondence} with $A=q$, $B=q^*$, and $c=m$.

\begin{lemma}[Exclusion of reciprocal symmetry]\label{lem:reciprocal}
There is no $s_0\in\Cs$ for which
\begin{equation}\label{eq:reciprocal}
 q(z)=q^*(s_0/z)\qquad(z\in\Cs).
\end{equation}
Consequently, no irreducible component of $Y$ has constant product
coordinate.
\end{lemma}

\begin{proof}
If \eqref{eq:reciprocal} held, the right side would tend to $m$ at
infinity. Hence $q$ would extend holomorphically at infinity and be
rational, with $q(\infty)=q(0)=m$.
The rational differential
\[
 \frac{q(z)}z\,dz=\frac{f'(z)}{f(z)}\,dz
\]
has residue $m$ at zero. At infinity, set $\zeta=1/z$. Then
\[
 \frac{q(z)}{z}\,dz=-\frac{q(1/\zeta)}{\zeta}\,d\zeta,
\]
so the residue is $-m$. If $a\ne0$ is a zero of $f$ of
multiplicity $n_a$, write $f(z)=(z-a)^{n_a}u(z)$ with
$u(a)\ne0$. Locally,
\[
 \frac{f'(z)}{f(z)}
 =\frac{n_a}{z-a}+\frac{u'(z)}{u(z)}.
\]
Thus every remaining pole has a positive integer residue $n_a$.
The residue theorem on the sphere gives
\[
 0=m-m+\sum_{a\ne0:\,f(a)=0}n_a.
\]
The sum is finite because the differential is rational. Since
all its summands are positive, there are no nonzero zeros of
$f$ and no other poles of $q$. The rational function $q$ is
therefore holomorphic on the sphere, hence constant $m$.
Finally, $q=m+zg'/g$ for $g=z^{-m}f$ gives $g'=0$, and
$f(z)=Cz^m$, a contradiction.

If an image component had $s=s_0$, a nonconstant source branch above
it would satisfy $w=s_0/z$. Its $z$-coordinate is nonconstant,
otherwise both coordinates would be constant. Its image contains an
open set of $z$-values, on which \eqref{eq:reciprocal} holds. The
meromorphic identity theorem gives \eqref{eq:reciprocal} on $\Cs$,
which has just been excluded.
\end{proof}

\begin{lemma}[Finiteness on each circle]\label{lem:finite-maxima}
A non-monomial entire function has finitely many maximum-modulus
points on each circle of positive radius.
\end{lemma}

\begin{proof}
The function $\theta\mapsto|f(re^{i\theta})|^2$ is real analytic.
If its maximum set were infinite, compactness would give an
accumulation point of the zeros of
$|f(re^{i\theta})|^2-M_f(r)^2$. The real-analytic identity theorem
would then make the modulus constant on the circle. In that case,
with $f^*(w)=\overline{f(\bar w)}$, the identity theorem gives
\[
 f(z)f^*(r^2/z)=M_f(r)^2\qquad(z\in\Cs).
\]
Indeed, the product is holomorphic on $\Cs$ and equals the positive
constant $M_f(r)^2$ on the circle, which has accumulation points in
that domain. The displayed identity excludes any nonzero zero of
$f$. Thus $g=z^{-m}f$ is a zero-free entire function and has constant
modulus $M_f(r)/r^m$ on the circle. The maximum-modulus principle
applied to both $g$ and $1/g$ gives, for $|z|<r$,
\[
 |g(z)|\le\frac{M_f(r)}{r^m},
 \qquad
 |g(z)|\ge\frac{M_f(r)}{r^m}.
\]
So $g$ is constant on the disk, hence everywhere, and $f$ is a
monomial. This contradiction proves that the original maximum
set is finite.
\end{proof}

\subsection{The almost-everywhere bound}
\begin{proposition}[The almost-everywhere bound]\label{prop:ae}
One has $\nu_f(r)\le2k$ for Lebesgue-almost every $r>0$.
\end{proposition}

\begin{proof}
Put $H(x)=\log M_f(e^x)$. Hadamard's three-circle theorem makes $H$
convex, hence locally Lipschitz and differentiable almost everywhere.
Fix a differentiability point $x$ and a global maximum
$z=e^{x+i\theta}$. Its value is nonzero, because $M_f(e^x)>0$.
Choose a holomorphic logarithm
$F(w)=\log f(e^w)$ near $w=x+i\theta$. Its derivatives satisfy
\begin{equation}\label{eq:logpartials}
 F'(w)=q(e^w),\qquad
 \frac{\partial}{\partial x}\Rea F(x+i\theta)=\Rea q(z),
 \qquad
 \frac{\partial}{\partial\theta}\Rea F(x+i\theta)=-\Ima q(z).
\end{equation}
Angular stationarity gives $\Ima q(z)=0$. On a neighborhood of $x$,
the nonnegative function
\[
 t\longmapsto H(t)-\log|f(e^{t+i\theta})|
\]
is differentiable at $x$ and has a local minimum zero there.
Its derivative at $x$ therefore vanishes. By
\eqref{eq:logpartials}, $H'(x)-\Rea q(z)=0$, so
\begin{equation}\label{eq:slope}
 q(z)=H'(x)\in\R.
\end{equation}
Thus each distinct global maximum gives a distinct pair $(z,\bar z)$
over
\begin{equation}\label{eq:maxtarget}
 (s,p)=(e^{2x},H'(x)).
\end{equation}
Here $z\bar z=e^{2x}$ and
$q^*(\bar z)=\overline{q(z)}=H'(x)$.

These targets stay away from the omitted value $p=m$ on each compact
$x$-interval. Indeed, write $f=z^m g$, where $g(0)\ne0$ and $g$ is
nonconstant. The function $J(x)=\log M_g(e^x)$ is convex and strictly
increasing. Strict increase follows from the maximum-modulus
principle applied to two concentric disks: if $M_g(r_1)=M_g(r_2)$
for $0<r_1<r_2$, a point on the inner circle attains the maximum
inside the outer disk, forcing $g$ to be constant. For a convex
function, slopes of successive secants are ordered, and the
increments $J(t+1)-J(t)$ are nondecreasing in $t$. Hence at a
differentiability point $x\in[a,b]$,
\[
 \begin{split}
 J(a)-J(a-1)
 &\le J(x)-J(x-1)\le J'(x)\\
 &\le J(x+1)-J(x)\le J(b+1)-J(b).
 \end{split}
\]
Strict increase makes the first increment positive. We obtain
\begin{equation}\label{eq:slopebounds}
 0<\delta_a:=J(a)-J(a-1)\le J'(x)
 \le C_b:=J(b+1)-J(b)<\infty.
\end{equation}
Since $H'=m+J'$, the targets in \eqref{eq:maxtarget} lie in the compact set
\[
 K_{a,b}=[e^{2a},e^{2b}]\times[m+\delta_a,m+C_b]\subset\Omega.
\]

The singular locus of the reduced curve $Y$ is locally finite, so
only finitely many singular image points lie in this compact set.
Their product coordinates exclude only finitely many radii. At every
remaining differentiability radius the target is smooth in the total
image, and Lemma~\ref{lem:degree}, Theorem~\ref{thm:correspondence}, and
Lemma~\ref{lem:reciprocal} give
\[
 \nu_f(e^x)
 \le\#\Phi^{-1}(e^{2x},H'(x))
 \le d_\Gamma\le2k.
\]
This excludes only a null set on each compact $x$-interval.
Exponentiation preserves null sets on such intervals, proving the
assertion for $r>0$.
\end{proof}

\subsection{Local constancy on either side of a radius}
The next elementary local argument turns the almost-everywhere
conclusion into the stated exceptional-set theorem. It allows angular
degeneracies and does not assume that $H$ is analytic through the
central radius.

\begin{lemma}[One-sided local constancy]\label{lem:onesided}
For every $r_0>0$, there exists $\delta>0$ such that $\nu_f$ is
constant on each of $(r_0-\delta,r_0)$ and $(r_0,r_0+\delta)$,
with $\delta<r_0$.
The two constants need not agree with each other or with $\nu_f(r_0)$.
\end{lemma}

\begin{proof}
By Lemma~\ref{lem:finite-maxima}, the central circle has finitely many
maxima. We first describe the stationary set near each of them, and
then show that these local descriptions cover all nearby maxima.

Fix one central maximum and choose a logarithmic coordinate $w$ near
it. Put $Q(w)=q(e^w)$ and let $w_0$ be the central point, with
$x_0=\Rea w_0=\log r_0$. The function $Q$ is holomorphic near
$w_0$ and nonconstant: constancy there would make $q$ constant
everywhere and $f$ a monomial. Angular stationarity gives
$p_0=Q(w_0)\in\R$. Factor
\[
 Q(w)-p_0=(w-w_0)^d h(w),\qquad d\ge1,\quad h(w_0)\ne0.
\]
Choosing a local holomorphic $d$th root of $h$ gives the conformal
coordinate $\chi(w)=(w-w_0)h(w)^{1/d}$, in which
\[
 Q(w)-p_0=\chi(w)^d,\qquad \chi(w_0)=0.
\]
Angular stationarity is $\Ima Q(w)=0$, so its local set consists of
the finitely many real-analytic half-arcs
\begin{equation}\label{eq:stationaryarcs}
 w_j(t)=\chi^{-1}(e^{i\pi j/d}t),
 \qquad t\ge0,\quad j=0,\ldots,2d-1.
\end{equation}
All half-arcs are restricted to sufficiently small $t$.

Choose disjoint angular windows $I_1,\ldots,I_{n_0}$ about the central
maxima so that their thin radial neighborhoods are contained in
the coordinate neighborhoods just constructed and avoid zeros of
$f$. Put $M_0=M_f(r_0)$. Compactness gives $\eta>0$ such that
\[
 |f(r_0e^{i\theta})|^2\le M_0^2-3\eta
 \qquad\left(\theta\notin\bigcup_{j=1}^{n_0} I_j\right).
\]
For $r$ sufficiently close to $r_0$, uniform continuity bounds the
left side with $r_0$ replaced by $r$ by $M_0^2-2\eta$ on this
complement. At any fixed central maximizing angle, the squared
modulus is greater than $M_0^2-\eta$. Every nearby global maximum
therefore lies in a window and belongs to one of the stationary
half-arcs in \eqref{eq:stationaryarcs}.

For each stationary half-arc, put $x_j(t)=\Rea w_j(t)$.
This function cannot be identically $x_0$.
Otherwise the stationary set would contain an open angular arc on
$|z|=r_0$. The modulus would be constant on that arc, hence on the
circle by real analyticity, contradicting
Lemma~\ref{lem:finite-maxima}. Therefore
\begin{equation}\label{eq:radialarc}
 x_j(t)-x_0=c_j t^{\ell_j}(1+t a_j(t)),
 \qquad c_j\in\R\setminus\{0\},\quad \ell_j\ge1.
\end{equation}
Here $a_j$ is real analytic near zero. Shrink the parameter interval
so that $1+t a_j(t)>0$, put $\sigma_j=\operatorname{sign}(c_j)$,
and define
\begin{equation}\label{eq:radialroot}
 \rho_j(t)=t\bigl[|c_j|(1+t a_j(t))\bigr]^{1/\ell_j}.
\end{equation}
The root is chosen positive for real $t$ near zero.
This is a real-analytic function with
$\rho_j'(0)=|c_j|^{1/\ell_j}>0$. The real-analytic inverse
function theorem gives its inverse $T_j=\rho_j^{-1}$, and
\[
 x_j(t)=x_0+\sigma_j\rho_j(t)^{\ell_j}.
\]
Thus the arc lies on the side $\sigma_j$ for small positive $t$
and intersects each sufficiently close radius on that side once.

Now include all half-arcs from all central windows, relabel them by
one finite index set, and choose $\ell_*$ to be a common multiple of
their finitely many exponents
$\ell_j$. Fix a side $\sigma\in\{-1,1\}$ and put
$x=x_0+\sigma v^{\ell_*}$ with $v>0$. Exactly the arcs with
$\sigma_j=\sigma$ contribute candidate points. Their parameters
and squared moduli are
\begin{equation}\label{eq:commonparameter}
 \begin{split}
 t_j(v)&=T_j\bigl(v^{\ell_*/\ell_j}\bigr),\\
 z_j(v)&=\exp\bigl(w_j(t_j(v))\bigr),\\
 b_j(v)&=|f(z_j(v))|^2.
 \end{split}
\end{equation}
Because $\ell_*/\ell_j$ is an integer, these functions are real
analytic at $v=0$. For small positive $v$ the points are
distinct: half-arcs in one window meet only at their center,
and the logarithmic coordinate is injective in each window;
different windows are disjoint. All nearby global maxima occur
among these finitely many candidates. Each side has at least one
candidate, since every nearby circle has a global maximum.

For any two applicable indices $i,j$, the real-analytic
difference $b_i(v)-b_j(v)$ either vanishes identically near zero
or has an expansion $a v^n+O(v^{n+1})$ with $a\ne0$.
Its sign is therefore fixed for all sufficiently small $v>0$.
There are only finitely many such differences, so one parameter
interval makes every comparison fixed. Consequently the set
of indices attaining the largest squared modulus is fixed
on that interval, and its cardinality is $\nu_f(e^x)$.
Do this separately on both sides and exponentiate. Reducing
$\delta$ so that $\delta<r_0$ proves the lemma.
\end{proof}

\begin{proof}[Proof of Theorem~\ref{thm:main}]
Fix $r_0>0$ and apply Lemma~\ref{lem:onesided}. Neither of its two
constant punctured-side values can exceed $2k$, by
Proposition~\ref{prop:ae}. Hence $r_0$ has a punctured neighborhood
containing no point of $E_f$. This holds at every positive radius,
so $E_f$ has no accumulation point in $(0,\infty)$ and is locally
finite. Sharpness is shown in Example~\ref{ex:sharp} below.
\end{proof}

\section{Sharpness and exceptional circles}\label{sec:examples}
\begin{example}[Sharpness]\label{ex:sharp}
For $k\ge1$, let
\[
 f_k(z)=\exp\!\left(z^k-\frac{z^{2k}}4\right).
\]
Its first coefficient gap is $k$, since $f_k(z)=1+z^k+O(z^{2k})$.
Set $\xi=z^k$, $|\xi|=R=r^k$, and
$T=\cos(\arg\xi)\in[-1,1]$.
Using $\cos(2\arg\xi)=2T^2-1$, complete the square:
\[
 \begin{split}
 \log|f_k(z)|
 &=\Rea\!\left(\xi-\frac{\xi^2}{4}\right)
 =RT-\frac{R^2}{4}(2T^2-1)\\
 &=\frac{R^2}{4}+\frac12-\frac12(RT-1)^2.
 \end{split}
\]
For $0<R\le1$, the square is minimized uniquely at $T=1$,
so the unique maximizing $\xi$-value is $\xi=R$. For $R>1$,
the minimum is zero at $T=1/R$, giving exactly the two values
$\xi=1\pm i\sqrt{R^2-1}$. At $R=1$ these two values coincide.
Every nonzero $\xi$ has exactly $k$ distinct preimages under
$z\mapsto z^k$, all of modulus $r$. Consequently,
\[
 \nu_{f_k}(r)=
 \begin{cases}
 k,&0<r\le1,\\
 2k,&r>1.
 \end{cases}
\]
Thus the constant $2k$ is attained on an entire tail.
It is also possible to see the degree bound directly in this
example. Here $q(z)=kz^k-\frac{k}{2}z^{2k}$ has real coefficients,
so $q^*=q$. Writing
$u=z^k$ and $v=w^k$, the correspondence factors as
\[
 q(z)-q(w)=k(u-v)\left(1-\frac{u+v}{2}\right).
\]
On the part $u+v=2$, one has $uv=s^k$ and
$p=k(u-u^2/2)=\frac{k}{2}s^k$.
For generic $s$, the equation $u(2-u)=s^k$ has two distinct
nonzero roots. Each root has $k$ choices of $z$, after which
$w=s/z$ is determined. The other factor $u=v$ meets this
image graph only when $u=v=1$ and $s^k=1$, so it contributes
no additional generic pairs. Hence this image component has
generic degree $2k$, exactly the bound of
Theorem~\ref{thm:correspondence}.
\end{example}

\begin{example}[Large multiplicity at an exceptional radius]\label{ex:exception}
Fix $N\ge3$ and $0<\tau<1/8$, and put
\[
 F(z)=\exp\!\left(z^N+\tau z(1-z^N)^4\right).
\]
Then $F(0)=1$, $F'(0)=\tau$, and the first coefficient gap is one.
Nevertheless, $\nu_F(1)=N$. To see this, let $d=|1-z^N|$ on $|z|=1$.
Since $|z^N|=1$, one has
$1-\Rea(z^N)=|1-z^N|^2/2=d^2/2$. Also,
$|\tau z(1-z^N)^4|=\tau d^4$ and $d\le2$, so
\[
 1-\Rea\bigl(z^N+\tau z(1-z^N)^4\bigr)
 \ge \frac{d^2}{2}-\tau d^4
 \ge\left(\frac12-4\tau\right)d^2.
\]
The coefficient $\frac12-4\tau$ is positive. Thus the real part
of the exponent is strictly below $1$ when $z^N\ne1$, and
equals $1$ at every $N$th root of unity. These are precisely
the $N$ maximum-modulus points. Differentiating the exponent
gives the logarithmic derivative
\[
 q(z)=Nz^N+\tau z(1-z^N)^4
             -4N\tau z^{N+1}(1-z^N)^3.
\]
Hence $F(\zeta)=e$ and $q(\zeta)=N$ for $\zeta^N=1$.
The fiber over $(s,p)=(1,N)$ therefore contains at least
$N$ distinct physical pairs $(\zeta,\bar\zeta)$.

This target point is singular in the total reduced image $Y$: at a
smooth point, Lemma~\ref{lem:degree} and
Theorem~\ref{thm:correspondence} would bound the fiber by two.
Several local image branches can meet at a singular point even when
they belong to the same global irreducible component. This is why
the proof removes singular image points before counting fibers.
\end{example}

Theorem~\ref{thm:main} bounds persistent multiplicity, not every
singular fiber. The examples above exhibit this distinction directly.
In particular, large values on isolated circles are compatible with
the theorem, whereas a lower bound tending to infinity on every
sufficiently large circle is not.

\subsection*{Disclosure of AI assistance}
OpenAI's ChatGPT and Codex (GPT-6) were used to develop the proof and
prepare the manuscript. The human author takes full responsibility for
the mathematical arguments and the final manuscript.

\begin{thebibliography}{99}

\bibitem{EvdoridouPardoSixsmith2021}
V.~Evdoridou, L.~Pardo-Sim\'on, and D.~J.~Sixsmith,
\emph{On a result of Hayman concerning the maximum modulus set},
Comput. Methods Funct. Theory \textbf{21} (2021), no.~4, 779--795.
\href{https://arxiv.org/abs/2012.07409}{arXiv:2012.07409}.

\bibitem{GlucksamPardo2024}
A.~Gl\"ucksam and L.~Pardo-Sim\'on,
\emph{An approximate solution to Erd\H{o}s' maximum modulus points problem},
J. Math. Anal. Appl. \textbf{531} (2024), no.~1, Paper No.~127768, 20~pp.
\href{https://arxiv.org/abs/2208.11154}{arXiv:2208.11154}.

\bibitem{Hayman1967}
W.~K.~Hayman, \emph{Research Problems in Function Theory},
The Athlone Press, University of London, London, 1967.

\bibitem{HerzogPiranian1968}
F.~Herzog and G.~Piranian,
\emph{The counting function for points of maximum modulus},
Entire Functions and Related Parts of Analysis,
Proc. Sympos. Pure Math., vol.~11,
Amer. Math. Soc., Providence, RI, 1968, pp.~240--243.

\bibitem{Lebl2026}
J.~Lebl, \emph{Tasty Bits of Several Complex Variables:
A Whirlwind Tour of the Subject}, version~4.4, May~31, 2026.
\url{https://www.jirka.org/scv/}.

\bibitem{Marchenko2012}
I.~I.~Marchenko,
\emph{On the maximum modulus points of entire and meromorphic functions
and a problem of Erd\H{o}s},
Mat. Stud. \textbf{38} (2012), no.~2, 212--215.
\href{https://doi.org/10.30970/ms.38.2.212-215}{doi:10.30970/ms.38.2.212-215}.

\bibitem{PardoSixsmith2026}
L.~Pardo-Sim\'on and D.~J.~Sixsmith,
\emph{A finite-order answer to a problem of Erd\H{o}s on maximum modulus points},
preprint, 2026.
\href{https://arxiv.org/abs/2607.09462v1}{arXiv:2607.09462v1}.

\end{thebibliography}
\end{document}
